\section{Extensions}
\label{sec:extensions}

In this section, we discuss some extensions which are all incorporated
in our proved mechanization. In particular, we cover sum types, replication and
recursion, recursive types, and the forwarding operator link. That is, all
features of GV \cite{conf/esop/LindleyM15} are covered.

\subsection{Sums}
\label{sec:extension:-sums}


\begin{figure}[tp]
  New values, terms, evaluation contexts \hfill \rocq{TSum}
  \begin{align*}
    \fFunc{v} &::= \dots \sepr \fFunc{\tInj{1}{v}} \sepr \fFunc{\tInj{2}{v}}
    & \text{injections}\\
    \fFunc{e} &::= \dots \sepr \fFunc{\tCase{e}{x}{e}{x}{e}} &
                                                                       \text{sum elimination}\\
    \fFunc{\fCtx{E}} & ::= \dots \sepr \fFunc{\tInj{1}{\fCtx{E}}} \sepr \fFunc{\tInj{2}{\fCtx{E}}}
  \end{align*}
  New labeled transitions \hfill \rocq{HeadCase1}
  \begin{mathpar}
    \begin{bussproof}[T-Act-Case]
      \bussAx{\fFunc{\tCase{(\tInj{i}{v})}{x_1}{e_1}{x_2}{e_2}} \trans{\beta} \fFunc{\tSub{e_i}{v}{x_i}}}
    \end{bussproof}
    % \and
    % \begin{bussproof}[T-Act-Case2]
    %   \bussAx{\fFunc{\tCase{\tInj{2}{v}}{\tLVA}{e}{\tLVA}{e}} \trans{\beta} \fFunc{\tSub{e'}{\tFixI{\tIVY}{e}}{\tIVX}}}
    % \end{bussproof}
  \end{mathpar}
  New typing rules for terms \hfill \rocq{typed_inj1}
  \begin{mathpar}
    \begin{bussproof}[E-Inj]
      \bussAs{\fType{\Sigma ; \Delta \vdash \fFunc{e} : T_i}}
      \bussU{\fType{\Sigma ; \Delta \vdash \fFunc{\tInj{i}{e}} : \lSum{T_1}{T_2}}}
    \end{bussproof}
    \and
    \begin{bussproof}[E-Case]
      \bussAs{\fType{\Sigma_1 ; \Delta_1 \vdash \fFunc{e} : \lSum{T_1}{T_2}}}
      \bussAsBelow{\fType{\Sigma_2 ; \Delta_2, \fFunc{x_1}:T_1 \vdash \fFunc{e_1} :U}}
      \bussAsBelow{\fType{\Sigma_2 ; \Delta_2, \fFunc{x_2}:T_2 \vdash \fFunc{e_2} :U}}
      \bussU{\fType{\Sigma_1 , \Sigma_2 ; \Delta_1 , \Delta_2 \vdash \fFunc{\tCase{e}{x_1}{e_1}{x_2}{e_2}} : U}}
    \end{bussproof}
  \end{mathpar}
  New case for the value logical relation \hfill \rocq{Vdef}
  \begin{align*}
    \lrVal{n}{\fType{\Sigma \vdash \lSum{T_1}{T_2}}}
  &= \set{ \fFunc{\tInj{1}{v}}  \mid
    \fFunc{v} \in \lrVal{n}{\fType{\Sigma \vdash T_1}}
    } \cup \set{ \fFunc{\tInj{2}{v}}  \mid
    \fFunc{v} \in \lrVal{n}{\fType{\Sigma \vdash T_2}}
    }
  \end{align*}
  New rule for primitive types \hfill \rocq{prim_sum}
  \begin{mathpar}
    \begin{bussproof}
      \bussAs{\primitive \fType{T}}
      \bussAs{\primitive \fType{U}}
      \bussB{\primitive \fType{\lSum{U}{V}}}
    \end{bussproof}
  \end{mathpar}
  \caption{Extensions for sum types}
  \Description{Additional values, terms and evaluation contexts for sums; one reduction axiom; two typing rules; one case of the value logical relation; and one rule for primitive types.}
  \label{fig:sum-types}
\end{figure}

Moving towards a realistic calculus, we wish to support branching
and choice in session types. It is well-known that branch and
choice can be simulated using higher-order sessions and sum types
\cite{DBLP:conf/ppdp/DardhaGS12,conf/esop/LindleyM15} by setting
\begin{align*}
  \fType{S_1 \ichoice S_2} &:= \fType{\sSend{(\dual{S_1} + \dual{S_2})}{\sTerm}}
  & \fType{S_1 \echoice S_2} &:= \fType{\sRecv{(S_1 + S_2)}{\sWait}}
\end{align*}
and encoding the operations to select and match on a channel as follows.
\begin{align*}
  \fFunc{\mathit{select}~\textit{left}~c} &:= \fFunc{\tFork (\lambda
                                   d. \tApp{\tClos}{(\tSend{\tPair{\tInj{1}{d}}{c}})})}
  &  \fFunc{\mathit{select}~\textit{right}~c} &:= \fFunc{\tFork (\lambda
                                       d. \tApp{\tClos}{(\tSend{\tPair{\tInj{2}{d}}{c}})})}
\end{align*}
\begin{align*}
  \fFunc{\mathit{match}~c~(x_1.e_1) (x_2.e_2)} &= \fFunc{\tLet[d][c]{\tRecv c}{\tClos{c}; \tCase{d}{x_1}{e_1}{x_2}{e_2}}}
\end{align*}
Directly supporting branch and choice would require additional labeled
reductions and hence additional cases in the definition of all logical
relations. To support the simulated branch and choice, it is sufficient to
implement sum types. To this end, we extend the syntax, labeled reduction,
typing, and the value logical relation in the usual way as
shown in \cref{fig:sum-types}. We prove the following compositionality lemmas
corresponding to the typing rules for sum types and integrate them in the proof
of the fundamental property. 

\begin{lemma}[Compositionality: \ruleLabel{inj} \rocq{sem_typed_inj1}]
  \label{lem:comp:inj}
  If
  $ \fType{\Sigma;\Delta \proves \fFunc{e} : T_i}$,
  then
  $ \fType{\Sigma;\Delta \proves \fFunc{\tInj{i}{e}} : \lSum{T_1}{T_2}}$,
  for $i\in\{1,2\}$.
\end{lemma}

\begin{lemma}[Compositionality: \ruleLabel{case} \rocq{sem_typed_case}]
  \label{lem:comp:case}
  If  $ \fType{\Sigma_1;\Delta_1 \proves \fFunc{e} : \lSum{T_1}{T_2}}$ and
  \begin{itemize}
  \item for all $ \fType{\Sigma;\Delta \proves \fFunc{v_1} : {T_1}}$ implies
    $ \fType{\Sigma,\Sigma_2;\Delta,\Delta_2 \proves \fFunc{e_1[v_1/x_1]} : U}$,
  \item for all $ \fType{\Sigma;\Delta \proves \fFunc{v_2} : {T_2}}$ implies
  $ \fType{\Sigma,\Sigma_2;\Delta,\Delta_2 \proves \fFunc{e_2[v_2/x_2]} : U}$,
  \end{itemize}
  then
  $ \fType{\Sigma_1,\Sigma_2;\Delta_1,\Delta_2 \proves
    \fFunc{\tCase{e}{x_1}{e_1}{x_2}{e_2}} : U}$. 
\end{lemma}

We exercise this extension with an example.

\begin{example}[Leader election]
The problem of leader election can be reduced to a process that makes
finite number of $\tau$ steps and then finishes with a $\fType{left}$
or $\fType{right}$ step on some channel of type $\fType{\fFunc{c} : \sTerm
  \ichoice \sTerm}$. To this end, we define the process $p_0 = (\set{\fFunc{cs},
  \fFunc{d}}, \mset{\fFunc{e_0}, \fFunc{e_1}, \fFunc{e_2}})$ which consists of
the following terms making use of the sum extension. We use standard syntactic
sugar for the encoding of the type $\fType{Bool} = \fType{\tUnit+\tUnit}$ and
the conditional $\fFunc{\mathit{if}}$.
%
\begin{align*}
  \fFunc{e_0} &=
                \fFunc{\tApp{\tClos}{(\tApp{\tSend}{\tPair{true}{\tApp{\tSend}{\tPair{false}{d}}}})}}
  \\
  \fFunc{e_1} &=
                \fFunc{\tLet[\tLVB][d]{\tApp{\tRecv}{d}}{\tLet*{\tApp{\tClos}{cs}}{\tIf{\tLVB}{\tApp{\tClos}{(\mathit{select}~\mathit{left}~c)}}{\tUnit}}}}
  \\
  \fFunc{e_2} &=
                \fFunc{\tLet[\tLVB][d]{\tApp{\tRecv}{d}}{\tLet*{\tApp{\tClos}{cs}}{\tLet*{\tApp{\tClos}{d}}{\tIf{\tLVB}{\tApp{\tClos}{(\mathit{select}~\mathit{right}~c)}}{\tUnit}}}}}
\end{align*}
Now we
set out to prove that $p_0 \in \lrProc{n}{ \fFunc{c}\fType{:  \sTerm
  \ichoice \sTerm} \vdash \fType{\tUnit}}$. This process has essentially
two different reduction sequences due to a race condition
(\rocq{leader_proc}). Either $\fFunc{e_0}$
communicates first with $\fFunc{e_1}$ and then with $\fFunc{e_2}$ or
vice versa. Each term, $\fFunc{e_1}$ and $\fFunc{e_2}$, receives a
boolean on channel $\fFunc{d}$, then synchronizes by closing
$\fFunc{cs}$, $\fFunc{e_2}$ closes $\fFunc{d}$, and we are left with
one term where the condition is $\fFunc{true}$ and another where it is
$\fFunc{false}$. The  $\fFunc{true}$-term executes a $\fFunc{select}$
and the $\fFunc{false}$-term does nothing.

The process $p_0$ is not reachable by a reduction sequence from a process of the
form $(\set{}, \mset{e})$, for some closed $e$, because
thread $\fFunc{e_0}$ shares $\fFunc{d}$ with both $\fFunc{e_1}$ and
$\fFunc{e_2}$. Term $\fFunc{e_1}$ is not typable
(cf. \cref{sec:syntacticTyping}) because it violates linearity for
variables/channels $\fFunc{d}$ and $\fFunc{c}$. Likewise, term
$\fFunc{e_2}$ violates linearity for channel $\fFunc{c}$.

See \cref{sec:reduct-lead-elect} in the appendix for the reduction sequence.
Composing $p_0$ over $\fFunc{c}$ with a syntactically typable client
yields a closed, deadlock-free process
(\rocq{leader_election_deadlock_free}).
\end{example}



\subsection{Recursion}
\label{sec:recursion}

\begin{figure}[tp]
  New values, terms, types \hfill \rocq{TBang}
  \begin{align*}
    \fFunc{v} &::= \dots \sepr \fFunc{\IExp{(\tIVX = e)}}
    & \text{recursively defined thunk}\\
    \fFunc{e} &::= \dots \sepr \fFunc{\tForce} & \text{thunk elimination}\\
    \fType{T} &::= \dots \sepr \fType{\tBox{T}}
  \end{align*}
  New labeled transition \hfill \rocq{HeadForce}
  \begin{mathpar}
    \begin{bussproof}[T-Act-Exp]
      \bussAx{\fFunc{\tApp{\tForce}{\IExp{(\tIVX = e)}}}
        \trans{\beta}
        \fFunc{\tSub{e}{\IExp{(\tIVX = e)}}{\tIVX}}}
    \end{bussproof}
  \end{mathpar}
  New typing rules for terms
  \hfill \rocq{typed_bang}
  \begin{mathpar}
    \begin{bussproof}[E-Exp]
      \bussAs{\fType{\tEmpty ; \tBox{\Delta} , \fFunc{\tIVX}:\tBox{T} \vdash \fFunc{e} : T}}
      \bussU{\fType{\tEmpty  ; \tBox{\Delta} \vdash \fFunc{\IExp{(\tIVX = e)}} : \tBox{T}}}
    \end{bussproof}
    \and
    \begin{bussproof}[E-Force]
      \bussAs{\fType{\tEmpty ; \tEmpty \vdash \fFunc{\tForce} : \lFunc{\tBox{T}}{T}}}
    \end{bussproof}
    \\
    \begin{bussproof}[E-Weak]
      \bussAs{\fType{\Sigma ; \Delta \vdash \fFunc{e} : U}}
      \bussU{ \fType{\Sigma ; \Delta, \tIVX : \tBox{T} \vdash \fFunc{e} : U}}
    \end{bussproof}
    \and
    \begin{bussproof}[E-Contract]
      \bussAs{\fType{\Sigma ; \Delta, \tIVX : \tBox{T}, \tIVY : \tBox{T} \vdash \fFunc{e} : U}}
      \bussU{ \fType{\Sigma ; \Delta, \tIVX : \tBox{T} \vdash \fFunc{\tSub{e}{\tIVX}{\tIVY}} : U}}
    \end{bussproof}
  \end{mathpar}
  \caption{Extending the term language and typing
    (\cref{fig:typeSys}) with recursive exponentials in the style of Lily}
  \Description{Additional values, terms and evaluation contexts for recursive exponentials; one reduction axiom; and three typing rules over a judgment carrying an extra intuitionistic environment.}
  \label{fig:lily:terms}
\end{figure}

We support term recursion by adapting ideas from Lily
\cite{DBLP:journals/entcs/BiermanPR00}, a linear lambda calculus with
recursion. Lily is based on a term calculus for linear logic with
exponentials where ``recursion is absorbed into terms of the
exponential type'' \cite{DBLP:journals/entcs/BiermanPR00}. In
consequence, we add a boxed type for duplicable values, an introduction form
for recursively defined values ``thunks'' in the style of Lily, and a $\fFunc{\tForce}$
operator as an eliminator for boxed types as shown in
\cref{fig:lily:terms}. Lily's eliminator is different, but equivalent.

The typing rules are as expected. Introducing a thunk requires an environment of
only boxed types, indicated by $\tBox{\Delta}$; forcing a thunk eliminates the
boxed type; boxed-typed variables are subject to weakening and contraction. The
mechanization implements the structural rules as part of environment splitting,
rather than explicit steps in a typing derivation.

Operationally, the thunk $\fFunc{\IExp{(\tIVX = e)}}$ is a recursively defined value and
the elimination operator exposes the body by unrolling the recursion.
% The reduction rule
% we propose here is proved to be a Kleene
% equivalence by \citet{DBLP:journals/entcs/BiermanPR00} with respect to
% their big-step semantics.
\citet{DBLP:journals/entcs/BiermanPR00} define some syntactic sugar for a fixed
point operator and non-recursive thunks:
\begin{align*}
  \fFunc{\tFixI{\tIVX}{e}} &:= \fFunc{\tApp{\tForce}{\IExp{(\tIVX = e)}}}
  & \fFunc{\IExp{e}} & := \fFunc{\IExp{(\tIVX = e)}} \quad \text{where}~\fFunc{\tIVX\notin fv (e)}
\end{align*}

We extend the value logical relation in a similar way as \citet{DBLP:journals/entcs/BiermanPR00}.
\begin{align*}
  \lrVal{n}{\fType{\tEmpty \vdash \tBox{T}}}
  &= \set{ \fFunc{\IExp{(\tIVX = e)}}  \mid
    \fc(\fFunc{e}) = \emptyset \land
    \fFunc{{\tFixI{\tIVX}{e}}} \in \lrTerm{n}{\fType{\tEmpty \vdash T}}
    }
\end{align*}
A replicable thunk owns no channels, which is what the side condition
$\fc(\fFunc{e}) = \emptyset$ records; exponentials are also admitted as
return types of a thread, so we add the rule
$\primitive \fType{T} \implies \primitive \fType{\tBox{T}}$
to the ones of \cref{sec:syntacticTyping:typeSys}.

For the additional cases of the fundamental theorem, we proved the
following compositionality lemmas, each corresponding to one of the typing rules
in \cref{fig:lily:terms}.

\begin{lemma}[Compositionality: \ruleLabel{exp} \rocq{sem_typed_bang}]
  \label{lem:comp:exp}
  If
  $ \fType{\tEmpty;\tBox{\Delta}, \tBox{T} \proves \fFunc{e} : T}$,
  then
  $ \fType{\tEmpty;\tBox{\Delta} \proves \fFunc{\IExp{(x=e)}} : T}$.
\end{lemma}

\begin{lemma}[Compositionality: \ruleLabel{force} \rocq{sem_typed_force}]
  \label{lem:comp:force}
  If
  $ \fType{\Sigma; \Delta \proves \fFunc{e} : \tBox{T}}$,
  then
  $ \fType{\Sigma; \Delta \proves \fFunc{\tApp{\tForce}{e}} : {T}}$.
\end{lemma}

\begin{lemma}[Compositionality: \ruleLabel{weak} \rocq{sem_typed_wk}]
  If
  $ \fType{\Sigma; \Delta \proves \fFunc{e} : U}$,
  then
  $ \fType{\Sigma; \Delta, x:\tBox{T} \proves \fFunc{e} : U}$.
\end{lemma}
\begin{lemma}[Compositionality: \ruleLabel{contract} % \rocq{sem_typed_contract}
  ]
  If
  $ \fType{\Sigma; \Delta, x:\tBox{T}, y:\tBox{T} \proves \fFunc{e} : U}$,\\
  then
  $ \fType{\Sigma; \Delta, x:\tBox{T} \proves \fFunc{\tSub{e}{x}{y}} : U}$.
\end{lemma}

This encoding of recursion also supports replication as in CP
\cite{conf/icfp/Wadler12,DBLP:journals/jfp/Wadler14}, the process calculus
companion to GV \cite{conf/esop/LindleyM15}. We implement a replicable
unrestricted session carrying type $\fType{S}$ as a function
of type $\fType{[S]} := \fType{\tBox{(\lFunc{\tUnit}{S})}}$ and
treat the new operators as macros as follows.
\begin{align*}
  \fFunc{repl} & \fType{{}:{\lFunc{\tBox{(\lFunc{S}{\tUnit})}}{[{\dual{S}}]}}}
  &
    \fFunc{req} & \fType{{}:{\lFunc{[S]}{S}}} \\
  \fFunc{repl} & :=
                 \fFunc{\tAbs{x}{\IExp{(\tAbs{\_}{\tApp{\tFork}{(\tApp{\tForce}{x})}})}}}
  & \fFunc{req} & := \fFunc{\tAbs{x}{\tApp{(\tApp{\tForce}{x})}{\tUnit}}}
\end{align*}

\mw{claude: likewise here --- examples/Factory.v mechanizes a
  replicated server: one factory value serving two independent
  clients, plus a discarded request, typed and deadlock free. It is
  the evidence that the encoding of Sect. 4.4 is usable end to end.
  Add it as an example?}
\pt{Yes, add this example here}

\subsection{Recursive Types}
\label{sec:extens-recurs-ii}

In the presence of recursive functions, it makes sense to extend
session types to be recursive, too. Specifically, we adopt the traditional approach of
equi-recursive session types, where a recursive type is always equal to its unfolding.
\begin{align*}
  \fType{S} &::= \dots \sepr \fType{\mu\alpha.S} \sepr \fType{\alpha}
  & \fType{\mu\alpha.S} &= \fType{\tSub{S}{\mu\alpha.S}{\alpha}}
\end{align*}
As equi-recursive types have contractive definitions
(\rocq{contractive_ty}), we can always
determine the top-level constructor of a type in finite time
(\rocq{unfoldS}). Hence,
all our prior definitions and proofs carry over without change. In the
mechanization, we insert appropriate unfolding functions that expose the
top-level constructor for inspection of the respective clause in the logical
relation.

\mw{claude: the mechanization has a worked example for this section
  that the paper does not: a recursive ping/pong over an equi-recursive
  session type, typed and deadlock free (examples/Pingpong.v). Add it
  as a short example, the way Sect. 4.1 has the leader election?}
\pt{this comment was for the previous subsection on recursion, but it makes more
sense to add it here where we have recursive types!}

\subsection{The Link Operator}
\label{sec:extens-link-oper}

\begin{figure}[tp]
  New constant, session label \hfill \rocq{ELink}
  \begin{align*}
    \fFunc{k} &::= \dots \sepr \fFunc{\tLink}
    & \text{forwarding operator}
    \\
    \sigma &::= \dots \sepr \lLink{c}{c}
             &\text{forwarding session label}
  \end{align*}
  New labeled transitions \hfill \rocq{HeadLink}
  \begin{mathpar}
    \begin{bussproof}[T-Act-Link]
      \bussAx{\fFunc{\tApp{\tLink}{\tPair{c}{d}}} \trans{\lLink{c}{d}} \fFunc{\tUnit}}
    \end{bussproof}
        \and
    \begin{bussproof}[P-Act-Link]
      \bussAs{\fFunc{e} \trans{\lLink{c}{d}} \fFunc{e'}}
      \bussU{\fProc{(C \cupdot \set{\fFunc{c}}, E \uplus \mset{\fFunc{e}})} \trans{\tau} \fProc{(C, \tSub{(E \uplus \mset{\fFunc{e'}})}{\fFunc{d}}{\fFunc{c}})}}
    \end{bussproof}
  \end{mathpar}
  Subjects and escaping channels of the new label:
  $ \fc(\lLink{c}{d}) = \set{\fFunc{c},\fFunc{d}} $, \quad
  $ \sent(\lLink{c}{d}) = \emptyset $

  New type-indexed family of constants: \hfill \rocq{KTLink}\\
  $ \fType{\fFunc{\tLink} : \lFunc{\tPair{S}{\dual{S}}}{\tUnit}} $

  New case for the term logical relation \hfill \rocq{Edef}
  \begin{enumerate}[start=6]
  \item  $\lambda = \lLink{c}{d}$ implies $\fFunc{c} \ne \fFunc{d}$ and $\fType{\Sigma} =
    \fType{\Sigma_1, \fFunc{c}:S, \fFunc{d}:\dual{S}}$ and $\fFunc{e'} \in \lrTerm{n'}{\fType{\Sigma_1 \vdash T}}$.
  \end{enumerate}

  New case for the process logical relation \hfill \rocq{Pdef}
  \begin{enumerate}[start=5]
  \item  $\pi = \lLink{c}{d}$ implies $\fFunc{c} \ne \fFunc{d}$ and $\fType{\Sigma} =
    \fType{\Sigma_1, \fFunc{c}:S, \fFunc{d}:\dual{S}}$ and $\fProc{p'} \in \lrProc{n'}{\fType{\Sigma_1 \vdash T}}$.
  \end{enumerate}

  The same two cases as obligation rules
  (\cref{fig:logrel-rules-term,fig:logrel-rules-proc})
  \begin{mathpar}
    \begin{bussproof}[O-Link]
      \bussAs{\fFunc{c} \ne \fFunc{d}}
      \bussAs{\fFunc{e'} \in \lrTerm{n}{\fType{\Sigma_1 \vdash T}}}
      \bussB{\obl{\fType{\Sigma_1, \fFunc{c}:S, \fFunc{d}:\dual{S}}}{n}{\fFunc{e} \trans{\lLink{c}{d}} \fFunc{e'} : \fType{T}}}
    \end{bussproof}
    \and
    \begin{bussproof}[Q-Link]
      \bussAs{\fFunc{c} \ne \fFunc{d}}
      \bussAs{\fProc{p'} \in \lrProc{n}{\fType{\Sigma_1 \vdash T}}}
      \bussB{\obl{\fType{\Sigma_1, \fFunc{c}:S, \fFunc{d}:\dual{S}}}{n}{\fProc{p} \trans{\lLink{c}{d}} \fProc{p'} : \fType{T}}}
    \end{bussproof}
  \end{mathpar}

  \caption{Extensions for forwarding (the link operator)}
  \Description{The new link constant and forwarding session label; the term and process reduction rules for forwarding; the subjects and escaping channels of the new label; the type of link; and the new case of the term and process logical relations, given both as a clause and as an inference rule.}
  \label{fig:link-operator}
\end{figure}

Session type systems based on the Curry-Howard correspondence for
classical linear logic (e.g.,
\cite{DBLP:journals/jfp/Wadler14,conf/esop/LindleyM15}) often include
a forwarding operator. This operator (called $\tLink$ by
\citet{conf/esop/LindleyM15}) takes two channel ends with dual
session types and connects them. This way, all messages on one channel
are forwarded to the other channel, in both directions.

\Cref{fig:link-operator} shows how to extend $\OurLang$ with
forwarding. We add a new operator $\fFunc{\tLink}$ and a forwarding label
$\lLink{\fFunc{c}}{\fFunc{d}}$ to indicate the request to forward between sessions
$\fFunc{c}$ and $\fFunc{d}$. As forwarding is symmetric the label
carries a \emph{set} of two channels, i.e., their order does not matter.
The two ends are \emph{not} required to be different: a self-forwarding
$\fFunc{\tApp{\tLink}{\tPair{c}{c}}}$ is a redex like any other and
emits the singleton label $\lLink{c}{c} = \set{\fFunc{c}}$.
The label is generated by a term
reduction. The corresponding process reduction requires that one of
the channel ends involved in forwarding is local. In this case, one
local channel ($\fFunc{c}$) is removed and the other channel is
substituted for $\fFunc{c}$. If no local channel is involved in the
label, the configuration as a whole takes a labeled reduction step
(cf. \TirName{P-Act-Session} in \cref{fig:proc-semantics}).
The new cases for the logical relations formalize the postconditition
of the forwarder: the two connected channel ends vanish from the
context as a consequence of the forwarding reduction.
Distinctness of the two ends is imposed by the logical relation, and
only there: a configuration offering a forwarding whose two ends
coincide names a single channel, which no instance of the case above
matches, so such a configuration is simply not related.
No typed term builds one either, since $\fFunc{\tLink}$ consumes a
session type together with its dual and linearity forbids one channel
supplying both.
On a restricted channel the substitution $\tSub{\cdot}{\fFunc{c}}{\fFunc{c}}$
is the identity, so \ruleLabel{P-Act-Link} degenerates to releasing the
name --- which is what a self-forwarding means: the thread held both
ends, so the session is spent.

\mw{claude: the mechanization carries a worked example for this
  section that the paper does not have --- Milner's cyclic scheduler
  with one agent replaced by a forwarder (Scheduler.v), closed,
  deadlock free, and untypable in Curry-Howard systems (it is
  Example~1 of Dardha and Gay). Add it as a subsection?}
Installing this extension, requires to update the proof of
\ruleLabel{E-Abs}
to cater for the new functional constant and
all other rules to cater for the new labeled reduction.

%%% Local Variables:
%%% mode: latex
%%% TeX-master: "main"
%%% End:
